\section{Future Work}
\label{sec:future-work}

\begin{itemize}

\item \textbf{Additional validated families.} The single highest-value
  extension (\S\ref{sec:limitations}, item~\ref{lim:n3}) is full manual
  reverse-engineering-backed validation of further families. A lighter,
  informal check was already run against a 4th, untuned family
  (AsyncRAT): a real CAPE detonation plus automatic resubmission of its
  dropped payloads confirmed the pipeline produces sane, non-crashing
  output end to end. This has no manual ground truth and is not treated
  as validation evidence, only an internal sanity check that the
  pipeline generalizes at all (documented in \code{README.md}); it does
  not substitute for the full, ground-truthed extension this item
  describes.

\item \textbf{Linux/ELF dynamic analysis for IoT botnet families.}
  Investigated concretely rather than dismissed by assumption, and
  deliberately deferred for three specific reasons, each independently
  verified: the pipeline's current ELF parser extracts basic structural
  metadata only, with no ATT\&CK detection rules wired to it at all,
  unlike the PE path's roughly 85 mapped techniques; CAPE's own
  documentation describes its Linux/ELF sandboxing support as an
  actively-developed, not production-grade, capability; and no
  ELF-family ground-truth sample exists yet against which any new
  Linux-specific detection rule could be verified, without which
  building rules would repeat exactly the overfitting risk
  \S\ref{sec:limitations} already flags for the Windows side. Given how
  disproportionately common ELF-based botnets are on IoT and embedded
  Linux devices relative to the malware-research literature's own
  attention to them, this is judged a genuinely valuable extension,
  just one that needs a validated ELF-family sample and a real dynamic
  analysis capability to exist first, not a rule set retrofitted onto
  static parsing alone.

\item \textbf{A statistically powered confidence-calibration study.}
  \S\ref{sec:confidence-calibration}'s finding that 14 of the 15
  audited discrepancies trace to single-source findings, with the one
  cross-source exception (RoningLoader's \technique{T1497}) attributable
  to a ground-truth completeness question rather than a miscalibrated
  confidence tier, is a directional signal from three samples, not a
  powered claim, and would benefit from a larger, purpose-built
  validation set.

\item \textbf{Runtime key-recovery via emulation.} The RC4 key that
  never exists as contiguous bytes in RoningLoader's DLL
  (\S\ref{sec:limitations}, item~\ref{lim:key-ceiling}) is, in principle, recoverable by
  emulating its assembly routine (e.g.\ with Unicorn) rather than any
  static byte-pattern search, a materially larger engineering effort
  than the current static config-recovery module, and scoped out of
  this thesis accordingly.

\end{itemize}
